\subsection{Pillar Correlation Analysis}
\label{app:pillar_correlations}

Table~\ref{tab:pillar-correlations} reports the full set of Pearson correlations between the inference-time knockout effect $\Delta_{\mathcal{B}}(v)$ (Eq.~\ref{eq:pillar-delta}) and nine candidate token-level predictors, measured over the $|\widetilde{\mathcal{R}}|=200$ screened Rock-Token candidates of \S\ref{sec:pillar-experiments} and reported separately for MATH-500 and IFEval. The figure in the main text (Figure~\ref{fig:pillar-vs-predictors}) visualizes a representative subset of six of these predictors on MATH-500; this table is the complete record, including IFEval and the three loss-based features (pre-/post-OPD KL and KL improvement). Across all eighteen $(predictor, benchmark)$ pairs, no correlation reaches $|r| > 0.075$ and no $p$-value falls below $0.30$. The largest observed magnitude on either benchmark is the IFEval correlation with KL improvement ($r{=}+0.071$, $p{=}0.32$), which is well within the noise band expected at $n{=}200$. We therefore cannot reject the null hypothesis $r=0$ for any predictor on either benchmark, supporting the claim in \S\ref{sec:pillar-vs-forking} that Pillarhood is statistically orthogonal to the entropy-, frequency-, and divergence-based importance signals used by prior work~\cite{wangbeyond,tip2026,yang2025token,wu2025mitigating}.

\begin{table}[h]
\centering
\caption{Pearson correlation between the per-token knockout effect $\Delta_{\mathcal{B}}(v)$ and nine candidate predictors over the $|\widetilde{\mathcal{R}}|=200$ screened Rock-Token candidates, reported for MATH-500 and IFEval. $p$-values are two-sided. No correlation reaches $|r|>0.075$ on either benchmark; no $p$-value falls below $0.30$.}
\label{tab:pillar-correlations}
\small
\begin{tabular}{lcccc}
\toprule
\textbf{Predictor} & \multicolumn{2}{c}{\textbf{MATH-500}} & \multicolumn{2}{c}{\textbf{IFEval}} \\
\cmidrule(lr){2-3}\cmidrule(lr){4-5}
 & $r$ & $p$ & $r$ & $p$ \\
\midrule
Frequency                  & $+0.015$ & $0.832$ & $+0.030$ & $0.677$ \\
Rock count                 & $+0.013$ & $0.855$ & $+0.023$ & $0.751$ \\
Rock rate                  & $-0.022$ & $0.760$ & $-0.047$ & $0.512$ \\
Pre-OPD KL                 & $-0.017$ & $0.808$ & $-0.015$ & $0.833$ \\
Post-OPD KL                & $+0.010$ & $0.883$ & $-0.048$ & $0.503$ \\
KL improvement             & $-0.062$ & $0.386$ & $+0.071$ & $0.320$ \\
Teacher entropy            & $+0.063$ & $0.378$ & $-0.036$ & $0.611$ \\
Pre-OPD student entropy    & $+0.059$ & $0.406$ & $-0.069$ & $0.331$ \\
Post-OPD student entropy   & $+0.056$ & $0.432$ & $-0.037$ & $0.604$ \\
\bottomrule
\end{tabular}
\end{table}

\subsection{Multiple-Testing Analysis for the Per-Token Pillar Criterion}
\label{app:pillar_multiple_testing}

The per-token Pillar criterion of \S\ref{sec:pillar-identification} applies a paired bootstrap test at $\alpha=0.05$ independently to each of the $|\widetilde{\mathcal{R}}|=200$ screened Rock-Token candidates, conjoined with the effect-size filter $|\Delta_\mathcal{B}(v)|\ge\varepsilon$ ($\varepsilon=0.01$). Because $\varepsilon$ is approximately at the per-token sampling-error scale, the conjunction does not appreciably reduce the per-test Type-I rate, and $200\!\times\!0.05 = 10$ rejections are expected under the global null. We report below how the per-token identifications behave under standard family-wise and false-discovery corrections, and the global sign-asymmetry test on which the population-level claim of \S\ref{sec:pillar-experiments} rests.

\paragraph{Per-token corrections.} Bonferroni control of the family-wise error at $\alpha=0.05$ requires $p<\alpha/200=2.5\times 10^{-4}$. The smallest observed $p$-values are $0.0038$ (MATH-500, ``\,certain'') and $0.0016$ (IFEval, ``-like''); neither passes Bonferroni, and consequently no token does. Benjamini--Hochberg control of the false discovery rate at $q\in\{0.05, 0.10, 0.20\}$ likewise produces zero rejections on either benchmark: the BH cutoff $p_{(k)}\le (k/n)\,q$ is not crossed even for the smallest-$p$ token, because $1/200 \times 0.20 = 10^{-3} < 0.0016$. The per-token Pillar identifications are therefore \emph{exploratory candidates}, not confirmatory rejections in the multiple-testing sense.



\paragraph{Global sign-asymmetry test.} The exploratory candidates do, however, exhibit a non-trivial structural property: the sign of the significant rejection is consistent across benchmarks. Of the seven MATH-500 candidates, all seven satisfy $\Delta_{\mathrm{MATH500}}(v)<0$ (Pillar direction); of the three IFEval candidates, all three satisfy $\Delta_{\mathrm{IFEval}}(v)<0$. Combined, the rejection sign-distribution is $10$ negative, $0$ positive. Under the global null in which every per-token rejection is a chance event, the rejection sign should be Bernoulli$(\nicefrac{1}{2})$, and the probability of observing all $10$ on one side is $2^{-10}\approx 9.8\times 10^{-4}$ (two-sided exact binomial $p = 1.95\times 10^{-3}$). On MATH-500 alone the analogous sign test is $7/0$ with two-sided $p=0.016$; on IFEval alone $3/0$ with $p=0.25$. The combined test is a single global test that does not inflate with the number of tokens screened, and it is the result on which \S\ref{sec:pillar-experiments} relies for the population-level claim that a non-empty subset of Rock Tokens is causally necessary at decode time.

\paragraph{Summary.} The reader is advised to read the per-token Pillar identifications of \S\ref{sec:pillar-experiments} (e.g.\ ``\,certain'', ``\,strategic'', ``\,Initialize'') as the strongest individual candidates from a population-level effect, rather than as $7$ independently-confirmed discoveries. The empirical contribution of this section that survives strict multiple-testing scrutiny is the sign asymmetry, which is incompatible with the multiplicity-noise null at $p\approx 0.002$.

\paragraph{Stable-core robustness check.} The screening pool of \S\ref{sec:pillar-identification} is the top-$K_{\text{screen}}=200$ count-ranked extension of $\mathcal{R}$, while the definitional Rock-Token set fixed by the stability/coverage analysis of \S\ref{ref:recalcitrant_tokens} is the top-$K=100$. The 10 Strong Pillars from the bootstrap distribute as follows over the count-rank index. On MATH-500 (ranks reported on the top-200 list): ``\,certain'' (16), ``\,Initialize'' (93), ``\,strategic'' (99), ``Do'' (109), ``\,programs'' (143), ``\,tech'' (166), ``\,arms'' (185). On IFEval: ``-like'' (35), ``\,tools'' (64), ``\,trade'' (69). Six of the ten lie within the $K=100$ core, including the smallest-$p$ token on each benchmark (``\,certain'' at $p=0.0038$ on MATH-500 and ``-like'' at $p=0.0016$ on IFEval). The remaining four lie in the rank-$109$-to-$185$ regime, which \S\ref{ref:recalcitrant_tokens} flags as below the Jaccard-stability sweet spot; we therefore caution that those four are at greater risk of being corpus-specific. The sign-asymmetry test of the previous paragraph is robust to this concern: restricted to the $K=100$ stable core, it remains $6/0$ negative-versus-positive, with two-sided exact-binomial $p = 2^{-6}\!\times\!2 = 3.1\times 10^{-2}$.
